\documentclass[12pt,a4paper]{article}
% ============================================================
%  Structural Persistence Series — Shared Preamble
% ============================================================
\usepackage{fontspec}
\setmainfont{Hiragino Mincho ProN}
\setsansfont{Hiragino Sans}
\setmonofont{Menlo}

\usepackage{iftex}
\ifXeTeX
  \XeTeXlinebreaklocale "ja"
  \XeTeXlinebreakskip=0pt plus 1pt minus 0.1pt
\fi

\usepackage[margin=22mm]{geometry}
\usepackage{amsmath,amssymb,amsthm}
\usepackage{xcolor}
\usepackage{graphicx}
\usepackage{hyperref}
\usepackage{booktabs}
\usepackage{fancyhdr}
% enumitem not available in this TeX installation

% --- Colours ------------------------------------------------
\definecolor{PaperAccent}{HTML}{1F4E79}
\definecolor{PaperGray}{HTML}{51606F}

\hypersetup{
  colorlinks=true,
  linkcolor=PaperAccent,
  urlcolor=PaperAccent,
  citecolor=PaperAccent,
  pdflang=ja
}
\urlstyle{same}

% --- Typography ---------------------------------------------
\setlength{\parindent}{1.5em}
\setlength{\parskip}{0pt}
\setlength{\emergencystretch}{3em}
\linespread{1.08}
\setlength{\abovedisplayskip}{8pt plus 2pt minus 4pt}
\setlength{\belowdisplayskip}{8pt plus 2pt minus 4pt}
\setlength{\abovedisplayshortskip}{6pt plus 2pt minus 2pt}
\setlength{\belowdisplayshortskip}{6pt plus 2pt minus 2pt}

% --- Section label names ------------------------------------
\renewcommand{\abstractname}{要旨}

% --- Theorem-like environments ------------------------------
\theoremstyle{plain}
\newtheorem{proposition}{命題}[section]
\newtheorem{lemma}[proposition]{補題}
\newtheorem{corollary}[proposition]{系}

\theoremstyle{definition}
\newtheorem{definition}{定義}[section]
\newtheorem{assumption}{条件}[section]

\theoremstyle{remark}
\newtheorem*{remark}{注}

% --- Running header -----------------------------------------
\pagestyle{fancy}
\fancyhf{}
\renewcommand{\headrulewidth}{0.4pt}
\fancyhead[L]{\small\textcolor{PaperGray}{\nouppercase{\leftmark}}}
\fancyhead[R]{\small\textcolor{PaperGray}{\thepage}}
\fancypagestyle{plain}{%
  \fancyhf{}
  \renewcommand{\headrulewidth}{0pt}
  \fancyfoot[C]{\small\textcolor{PaperGray}{\thepage}}
}

% --- Title block macro --------------------------------------
\newcommand{\PaperTitleBlock}[3]{%
  % #1 = title, #2 = subtitle, #3 = date
  \thispagestyle{plain}%
  \begin{center}
  {\LARGE\bfseries #1\par}
  \vspace{0.4em}
  {\normalsize #2\par}
  \vspace{1.0em}
  {\large Akihito Sunagawa\par}
  \vspace{0.3em}
  {\normalsize #3\par}
  \end{center}
  \vspace{1.6em}
}

% Legacy 2-arg version (backward compat)
\newcommand{\PaperTitleBlockSimple}[2]{%
  \PaperTitleBlock{#1}{}{#2}
}


\begin{document}

\PaperTitleBlock
  {構造持続写像の標準手順}
  {— ドメイン横断適用のための手順 —}
  {2026年4月12日}

\begin{abstract}
本補論は、構造持続の最小形式

\[
S = Me^{-L}
\]

を個別ドメインへ写像するための標準手順を与える。目的は、任意の対象について、それが構造維持系として扱えるか、崩壊を維持可能性の喪失として定義できるか、そしてそのドメインについて、定理、境界、または実証のいずれの強さで構造持続の数学に載せられるかを、同じ手順で判定できるようにすることである。

本手順は、対象構造、維持条件、有効選択肢空間、制約集合、累積構造損失 L、余力 μ、初期有効選択肢多様性 N\_eff\^{}(0) を順に定義し、Route A/B/C のいずれかへ位置づける。さらに、反証条件の事前固定、基準モデルとの比較、予測頑健性テスト、介入設計までを一つの流れとして与える。これにより、ドメインごとに異なる語で語られてきた崩壊や劣化を、共通の骨格のもとで比較・検証するための標準手順を与える。
\end{abstract}

\section{目的}\label{ux76eeux7684}

本補論は、任意のドメインについて、

\begin{itemize}
\tightlist
\item
  その対象が「構造維持系」として扱えるか
\item
  崩壊を「維持可能性の喪失」として定義できるか
\item
  構造持続の数学に定理的、境界的、または実証的に載せられるか
\end{itemize}

を判定し、実験・定式化・介入設計まで一貫して行うための標準手順を定めるものである。

本手順が扱うのは、存在一般ではなく、ある対象構造がその関係・機能・同一性を保ったまま持続できるかどうかである。

この手順の基本発想は、構造持続の最小形式で与えられた

\[
S = Me^{-L}
\]

を、SAT や LLM 以外のドメインにも写像できるかを系統的に判定することにある。

運用上は、M をさらに分解した形

\[
S = N_{eff}^{(0)} \times (\mu/\mu_{c}) \times e^{-L}
\]

を用いる。ここで \(M = N_{eff}^{(0)}\) × (μ/μ\_c) である。\(\mu_{c}\) は、そのドメインで持続が急速に失われ始める臨界余力を表す。この分解は最小形式の拡張ではなく、M の内部構造を操作可能にするための実用的な展開である。ここで N\_eff\^{}(0) が表すのは、制約適用後の現在の選択肢数ではなく、縮小が始まる前にその系が持っていた初期的または外生的な有効選択肢多様性である。現在の有効選択肢の縮小は e\^{}\{-L\} 側に含まれる。したがって、μ と N\_eff\^{}(0) の分離が困難なドメインでは、M を一体のまま扱ってよい。

\section{基本思想}\label{ux57faux672cux601dux60f3}

\subsection{中核仮説}\label{ux4e2dux6838ux4eeeux8aac}

任意の構造維持系は、次のように捉えられる。

\begin{itemize}
\tightlist
\item
  系には維持すべき機能がある
\item
  その機能を維持するための有効な行動・状態・経路の集合がある
\item
  制約や矛盾の蓄積は、その集合を削る
\item
  回復余地や選択肢が閾値以下になると崩壊する
\end{itemize}

\subsection{崩壊の一般定義}\label{ux5d29ux58caux306eux4e00ux822cux5b9aux7fa9}

崩壊とは、単なる停止ではなく、

「定義された維持機能を支える有効選択肢空間が、閾値以下に縮退すること」

である。

\subsection{L の一般定義}\label{l-ux306eux4e00ux822cux5b9aux7fa9}

L とは、単なる矛盾数ではなく、

「維持可能な有効選択肢を削る累積構造損失」

である。最小形式の累積損失 L と同一の量であり、情報量として見れば構造損失情報量にあたる。

\subsection{μ の一般定義}\label{ux3bc-ux306eux4e00ux822cux5b9aux7fa9}

μ とは、

「修復、再編、判断、行動、推論、適応に使える有効余力」

である。

\subsection{N\_eff\^{}(0) の一般定義}\label{n_eff0-ux306eux4e00ux822cux5b9aux7fa9}

N\_eff\^{}(0) とは、

「縮小が始まる前に、その系が持っていた独立な有効選択肢の初期的または外生的な多様性」

である。

μ と N\_eff\^{}(0) は最小形式の M（有効維持資源）を構成する二つの側面であり、μ は余力の絶対量、N\_eff\^{}(0) は初期または外生的な選択肢多様性に対応する。現在の有効選択肢の減少そのものは e\^{}\{-L\} に含まれる。両者を分離できるドメインでは分けて扱い、分離困難なドメインでは M を一体のまま扱う。

ただし、学習・進化・組織再編のように、選択肢多様性そのものが時間とともに再生・拡張される系もある。この場合、その再生が回復・再編能力と切り分けにくければ μ 側に含めてよく、独立に追跡できるときに限って将来的には時間依存の N\_eff(t) として拡張しうる。本手順では、まず基準形として N\_eff\^{}(0) を用いる。

\subsection{量の切り分け}\label{ux91cfux306eux5207ux308aux5206ux3051}

本補論では、量の種類を次のように区別する。

\begin{itemize}
\item
  \(R_{k}\): 有効選択肢空間の総残存比。正確に測れる場合には \$R\_\{k\} = \(m(\)A\_\{k\}\() / m(A_{k}^0)\) = e\^{}\{-L\_\{k\}\}\$ と書かれる。
\item
  \(M_{k}\): 有効維持資源。余力や初期的・外生的な多様性を含む、構造持続を支える資源項である。
\item
  \(S_{k}\): 構造持続ポテンシャル。 \$S\_\{k\} = \(M_{k}\) R\_\{k\}\$ で与えられる。
\item
  \(P_{k}\): Route B で用いる、確率的または弱依存の記述のもとでの実現残存可能性。役割としては Route A の \(R_{k}\) に対応するが、参照モデルに対する境界として表すときにこの記号を用いる。
\end{itemize}

したがって、L は「どれだけ削られたか」、R は「どれだけ残っているか」、M は「その縮小が始まる前からどれだけの余力と多様性を持っていたか」、S は「その両者を合わせた構造持続ポテンシャル」を表す。

\section{適用対象}\label{ux9069ux7528ux5bfeux8c61}

\subsection{適用可能な系}\label{ux9069ux7528ux53efux80fdux306aux7cfb}

\begin{itemize}
\tightlist
\item
  維持すべき機能が定義可能
\item
  崩壊が観測可能または操作的に定義可能
\item
  行動、状態、推論、意思決定などの有効選択肢を定義可能
\item
  制約や矛盾の蓄積が、その有効選択肢を減らすと考えられる
\end{itemize}

\subsection{典型例}\label{ux5178ux578bux4f8b}

\begin{itemize}
\tightlist
\item
  LLM
\item
  継続学習システム
\item
  自律エージェント
\item
  企業・組織
\item
  生態系
\item
  医療運用系
\item
  サプライチェーン
\item
  個体認知・意思決定系
\end{itemize}

\subsection{適用が難しい対象}\label{ux9069ux7528ux304cux96e3ux3057ux3044ux5bfeux8c61}

逆に、次のような対象は本手順にそのまま載せにくい。

\begin{itemize}
\tightlist
\item
  維持すべき機能そのものが定義できない対象
\item
  崩壊条件を観測的または操作的に書けない対象
\item
  有効選択肢空間を置く意味がほとんどない対象
\item
  制約蓄積よりも、生成・揺らぎ・再編が支配的で、削減方向を安定に定義しにくい対象
\item
  対象構造を後から都合よくずらせてしまう対象
\end{itemize}

例として、理想的な i.i.d. ランダム列、暗号学的乱数源の出力列そのもの、あるいは評価基準を与えない自由生成の列は、本手順にそのまま載せにくい。これらは、維持すべき機能、崩壊条件、有効選択肢空間のいずれも安定に定義しにくく、制約蓄積による持続失敗として記述するより、独立試行や生成過程として扱う方が自然だからである。

とくに注意すべきなのは、操作的に定義できない隠れ本質を事後的に持ち込むことである。この種の仮定は、対象構造をいくらでも後付けで再定義できてしまうため、反証可能性を失わせる。本手順が扱うのは、そのような形而上学的な本質ではなく、少なくとも維持条件・崩壊条件・有効選択肢のいずれかを操作的に固定できる構造である。

ただし、現時点で直接観測できない内部機構があること自体は、直ちに本手順の対象外を意味しない。重要なのは、その内部機構を仮定しなくても、維持条件、崩壊条件、介入効果のいずれかが操作的に記述できるかどうかである。したがって、「まだ完全には見えていないが、代理指標と介入設計を通じて検証可能な構造」は Route B または Route C の候補になりうるが、「見えない本質」そのものを対象構造に据えることは許されない。

\section{ドメイン特性の分類}\label{ux30c9ux30e1ux30a4ux30f3ux7279ux6027ux306eux5206ux985e}

対象ドメインを、有効選択肢空間の明示性と制約削減の測定可能性に基づいて分類する。この分類は手順 8 のルート判定の入力となる。

定理化寄りのドメイン: - 維持条件が外部的に明示定義できる - 有効選択肢空間が比較的明示的 - 制約による削減率が直接または近似的に計算できる - → Route A または Route B に進みやすい

例: SAT、明示制約型 LLM タスク、一部の形式推論環境

代理指標実証寄りのドメイン: - 維持条件は操作的に定義可能 - 有効選択肢空間は暗黙的または高次元 - 制約削減は直接測定困難で、代理指標が必要 - → Route B または Route C に進みやすい

例: 企業、生態系、一般的な LLM 会話、継続学習システム、組織運営

\section{写像手順}\label{ux5199ux50cfux624bux9806}

手順 0: 対象ドメインと対象構造の固定

対象を一つに絞り、対象構造を明記する。対象構造とは、「そのドメインにおいて何の構造持続を問題にするか」を指す。ここでいう構造は、一般的な形そのものではなく、その系で維持を問題にする関係・機能・同一性を指す。ドメインと対象構造は異なる。たとえばドメインが「LLM 対話系」でも、対象構造は「論理一貫性を保った推論構造」かもしれないし、「事実再現の整合構造」かもしれない。企業なら「継続的意思決定構造」、生態系なら「生態機能を支える相互作用構造」のように具体化される。同一ドメインに対して異なる対象構造を設定すれば、適用結果は変わりうる。

対象構造は、観測指標や代理指標を見る前に固定しなければならない。そうでなければ、うまく当てはまる対象構造が見つかるまで定義をずらし続けることができてしまう。

このとき固定されるべきなのは、このドメインで何を維持できればその構造が存続しているとみなすか、という操作的な対象構造である。対象構造は完全に可視化されている必要はないが、少なくとも崩壊条件や介入効果と結びつく形で固定されていなければならない。

出力: - ドメイン名 - 対象構造 - 対象スコープ - 観測期間または観測単位

手順 1: 維持条件の定義

対象ドメインで維持すべき機能を一つ定義する。

要件: - 一つでよい。最初から総合知能や総合健全性を定義しなくてよい - 操作的に観測可能であることが望ましい

例: - LLM: 論理一貫性維持 - 企業: 継続的意思決定能力維持 - 生態系: 生態機能維持 - 継続学習: 旧知識を壊さず新知識を統合する能力維持

出力: 維持条件 \(G_{k}\)

手順 2: 崩壊の定義

崩壊を、その維持条件 \(G_{k}\) が失われる条件として定義する。

推奨定義: 「維持条件を満たすための有効選択肢が閾値以下になること」

出力: 崩壊条件 \(C_{k}\)

手順 2.5: 反証条件の事前設定

\textbf{崩壊定義の直後、代理指標探索の前に}、以下を先に固定する。

\begin{itemize}
\tightlist
\item
  何が観測されたら「この理論はこのドメインに当てはまらない」と判断するか
\item
  何が観測されたら「選んだ代理指標が不適切」と判断するか
\end{itemize}

この区別が重要である。理論の棄却と代理指標の棄却は別の判定であり、代理指標が外れたことは直ちに理論の棄却を意味しない。ただし、すべての合理的な代理指標が外れた場合は、理論の適用可能性を再検討する。

反証条件を結果の後に書くことは許されない。

加えて、代理指標探索の上限を事前に固定する。「N 個の合理的な代理指標を試して全て外れたら、このドメインへの理論の適用可能性を再検討する」という事前コミットメントがなければ、「この代理指標は合理的とは言えなかった」という判定を繰り返すことで理論を永遠に守り続けることができてしまう。

出力: - 反証条件（理論レベル + 代理指標レベル） - 代理指標探索予算（試行する代理指標の上限数または探索期限）

手順 3: 有効選択肢空間の定義

維持条件 \(G_{k}\) を支える有効選択肢空間 \(A_{k}\) を定義する。

候補: - 有効な軌跡 - 実行可能な行動分岐 - 整合的な推論経路 - 回復可能な状態遷移 - 方針整合的な意思決定経路

要件: - 明示的でも暗黙的でもよい - 必ずしも全列挙可能でなくてよい - ただし何を''有効''と呼ぶかは定義する

出力: 有効選択肢空間 \(A_{k}\)

手順 4: 制約の定義

有効選択肢を削る制約 \(C_{i}\) を定義する。

候補: - 矛盾 - 上書き - 依存衝突 - 資源律速 - 制度的不整合 - 指示衝突 - 時間的不整合 - 範囲づけなし更新

重要: 制約は「悪い現象」ではなく、有効選択肢を削る作用として定義する。

出力: 制約集合 \{C\_i\}

手順 4.5: 基準モデルの設定

本手順の指数型が何に対して優位かを明確にするため、各ドメインで最も単純な代替モデルを先に設定する。

\begin{itemize}
\tightlist
\item
  「L を使わず、最も単純なモデルで崩壊を予測するとしたら何か」
\item
  「本手順はその基準モデルに対して何を追加で予測するか」
\end{itemize}

例: - LLM: 「矛盾の数が閾値を超えたら崩壊」（線形カウントモデル） - 企業: 「赤字が N 期連続したら崩壊」（単純閾値モデル）

本手順が基準モデルに勝てない場合、その事実自体が理論を前進させる情報になる。指数型の加法的累積構造が、単純なカウントや線形モデルに対して予測精度で優位であることを示せなければ、理論の実質的な貢献は「当たり前の観察に複雑な言語を与えただけ」になりうる。

基準モデルに対する優位性は、少なくとも次のいずれか一つで示されなければならない。 - held-out データでの予測優位 - 制約の質に関する順序予測の的中 - 介入順位予測の的中

出力: 基準モデル、追加予測

手順 5: 削減率または代理指標の定義

各 \(C_{i}\) が \(A_{k}\) をどれだけ削るかを定義する。

指数形式そのものは定義から従う。経験的に検証されるのは、どの制約型がより強く有効選択肢を削るか、その順序予測が保たれるか、そして高次元ドメインで観測される代理指標が真の累積構造損失の十分よい代理となるかである。

理想形: \[
\begin{aligned}
r_{i} &= |A_{k}^{(i)}| / |A_{k}^{(i-1)}| \\
L_{k} &= \sum[-\ln(r_{i})] \\
R_{k} &= e^{-L_{k}} 
\end{aligned}
\]

ここでの \(L_{k}\) は、真の累積構造損失である。

代理指標形（直接 \(r_{i}\) が測れない場合）:

直接 \(r_{i}\) が測れない場合は、真の \(L_{k}\) そのものではなく、代理累積構造損失 L̂\_k を置く。L̂\_k に求めるのは真の \(L_{k}\) との完全一致ではなく、少なくとも次のいずれかを再現することである。

\begin{itemize}
\tightlist
\item
  維持指標に対する単調関係
\item
  閾値予測
\item
  介入順位予測
\item
  基準モデルに対する追加予測力
\end{itemize}

代理指標の例: - 矛盾密度 - 未解決衝突数 - 方針衝突数 - 検索不一致率 - 意思決定膠着率 - 経路無効化率

出力: 真の \(L_{k}\) または代理累積構造損失 L̂\_k

手順 6: μ 候補の定義

系が修復・再編・行動・推論に使える余力 μ を定義する。

必要に応じて、臨界余力 \(\mu_{c}\) も定義する。\(\mu_{c}\) は、そのドメインで維持条件が急速に失われ始める余力の基準値である。

候補: - 計算余力 - 注意余力 - キャッシュ余力 - 資金余力 - 冗長性 - バッファ - 回復帯域

出力: 余力変数 μ\_k（必要に応じて \(\mu_{c}\),k）

手順 7: N\_eff\^{}(0) 候補の定義

系が、縮小が始まる前に利用可能だった独立な選択肢の多様性を定義する。

候補: - 代替政策の初期数 - 推論経路の初期多様性 - 冗長モジュールの初期数 - サプライ代替の初期数 - 記憶表現の初期多様性

ここで定義する N\_eff\^{}(0) は、現在残っている選択肢数ではない。現在の有効選択肢数の減少は L によって表される。μ と N\_eff\^{}(0) の分離が困難な場合は、手順 6-7 を統合し M として一体で扱う。

出力: 初期有効選択肢多様性 N\_eff\^{}(0)\_k（または統合 \(M_{k}\)）

手順 8: 数学ルート判定

次の3ルートのいずれかに分類する。ルート分類はドメインそのものの分類ではなく、そのドメインについて本稿でどの強さの主張が可能かを示す分類である。ルートの複数性は逃げ道ではなく、主張強度を下げる仕組みである。Route A では等式、Route B では境界、Route C では経験的傾向が主張される。ルートが下がるほど主張は弱くなるが、反証条件と基準モデル比較が機能する限り、理論の適用可能性は検証可能なまま保たれる。

Route A: 定理ルート

条件: - \(A_{k}\) が定義可能 - 各段階の残存比 \(r_{i}\) が定義可能 - 累積構造損失 \(L_{k}\) = Σ{[}-ln(\(r_{i}\)){]} が定義可能 - 共通測度のもとで残存集合の比較ができる

目標: \textbar A\_k（残存）\textbar{} = \textbar A\_k\^{}0\textbar{} e\^{}\{-L\_k\}

Route B: 境界ルート

条件: - 等式としての Route A は置けない - ただし、参照累積損失 \(L_{ref}\) からの相対誤差が有界

目標: \[
e^{-L_{ref}(1+\rho)} \le P_{k} \le e^{-L_{ref}(1-\rho)}
\]

既存理論の弱依存境界を用いる。ρ は段階損失間の依存が参照累積損失に及ぼす相対誤差の上界である。ここで \(P_{k}\) は、弱依存や確率モデルのもとで記述された実現残存可能性を表す。

Route C: 実証ルート

条件: - 厳密な定理はまだ難しい - ただし代理指標による予測・介入・再現が可能

目標: - 単調性 - 閾値性 - 介入効果 - 予測力 - \textbf{基準モデルに対する優位性}

出力: ルート分類

手順 9: 代理指標探索フェーズ

Route C（および Route B で代理指標を併用する場合）で必須。

手順: 1. 候補の代理指標を複数立てる 2. 維持指標との相関を見る 3. 単調性を見る 4. 閾値性を見る 5. 外れた代理指標を捨てる \# \textbf{基準モデルと比較し、代理指標が追加の予測力を持つか検証する}

出力: - 候補代理指標集合 - 棄却代理指標集合 - 生存代理指標集合

複数の代理指標を試行した場合、生存した代理指標の予測力は新規データで再検証する。

手順 10: 検証フェーズ

生き残った代理指標を使い、以下を検証する。

\begin{itemize}
\tightlist
\item
  他データで再現するか
\item
  他モデルで再現するか
\item
  維持条件の定義を変えても効くか
\item
  介入と連動するか
\item
  未知データで予測できるか
\end{itemize}

出力: 検証済み操作的 L̂\_k

手順 10.5: 予測頑健性テスト

高次元ドメインでは、有効選択肢空間 \(A_{k}\) それ自体を直接観測できないことが多い。この場合、空間定義の妥当性は、空間そのものの直接比較ではなく、その定義から導かれる予測の頑健性によって評価される。

具体的には、以下の定義パラメータを複数の合理的な選択肢で変えたとき、主要予測が保たれるかを調べる。

\((a)\) 有効選択肢の粒度 例: 推論経路を「前提→規則→結論」の3段階で数えるか、 各推論ステップを独立に数えるか。 \((b)\) 測度の取り方 例: 候補経路を一様に数えるか、確率重みをつけるか。 \((c)\) 制約の分類法 例: 制約を5分類にするか、3分類にするか、2分類にするか。

検証対象となる予測: - 制約の質の順序予測（どの制約型がより強く削るか） - 代理指標と維持指標の単調関係 - 介入効果の方向と相対的大きさ - 基準モデルに対する追加予測力

判定基準: これらの予測が、(a)(b)(c) の合理的な変更に対して鈍感であれば、骨格は頑健であるとみなす。逆に、粒度や測度や分類の選び方を少し変えるだけで予測が反転するならば、理論の実質は骨格ではなく定義選択に依存しており、そのドメインでの適用は未成熟であると判断する。

この手順は、有効選択肢空間が明示的に定まるドメイン（例: SAT）では不要である。空間の定義がほぼ一意に決まるため、感度分析は自明に通過する。本手順が必要になるのは、有効選択肢空間が理論構成物として置かれ、定義に自由度が残るドメインにおいてである。

出力: 各定義パラメータの変更に対する予測の安定性報告

手順 11: 介入設計

理論を因果的に強くするため、以下の介入を設計する。

L を減らす介入: - 矛盾解消 - 時間的範囲づけ - 記憶標識 - 巻き戻し - 剪定 - 衝突仲裁

μ を増やす介入: - バッファ拡大 - 計算予算追加 - 冗長性確保 - 資源余裕投入

N\_eff\^{}(0) を確保・拡張する介入: - 代替経路確保 - 表現多様性維持 - モジュール化 - 方針分岐保持

検証: 介入後に維持指標が改善するかを見る。

\begin{enumerate}
\def\labelenumi{\arabic{enumi}.}
\setcounter{enumi}{5}
\tightlist
\item
  出力フォーマット
\end{enumerate}

本手順を適用した各ドメインは、以下のフォーマットで整理する。

ドメイン構造写像シート:

{\def\LTcaptype{table} % do not increment counter
\begin{longtable}[]{@{}
  >{\raggedright\arraybackslash}p{(\linewidth - 2\tabcolsep) * \real{0.5000}}
  >{\raggedright\arraybackslash}p{(\linewidth - 2\tabcolsep) * \real{0.5000}}@{}}
\toprule\noalign{}
\begin{minipage}[b]{\linewidth}\raggedright
項目
\end{minipage} & \begin{minipage}[b]{\linewidth}\raggedright
内容
\end{minipage} \\
\midrule\noalign{}
\endhead
\bottomrule\noalign{}
\endlastfoot
ドメイン & 対象ドメイン名 \\
対象構造 & 何の構造持続を問題にするか \\
維持条件 & 維持すべき機能 \\
崩壊条件 & 崩壊の定義 \\
反証条件 & 理論レベル + 代理指標レベル \\
代理指標探索予算 & 試行する代理指標の上限数または探索期限 \\
基準モデル & 最も単純な代替予測モデル \\
追加予測 & 基準モデルに対する本手順の追加予測 \\
有効選択肢空間 & 定義 \\
制約集合 & 定義 \\
L / L̂ 定義 & 真の累積構造損失 L または代理累積構造損失 L̂ の定義 \\
μ 定義 & 余力定義 \\
N\_eff\^{}(0) 定義 & 初期有効選択肢多様性定義（分離困難なら M として統合） \\
ルート & 定理 / 境界 / 実証 \\
観測指標 & 定義 \\
介入候補 & 定義 \\
検証計画 & 定義 \\
\end{longtable}
}

\section{判定基準}\label{ux5224ux5b9aux57faux6e96}

この標準手順がそのドメインで「うまく働く」とみなす条件は次の通り。

最低条件: - 維持条件が操作的に定義できる - 崩壊が観測可能 - 制約候補が定義できる

良条件: - 代理累積構造損失 L̂ が維持指標を単調に予測する - 介入が維持指標の改善と連動する - \textbf{基準モデルに対して追加の予測力がある}

強条件: - 有効選択肢空間と削減率が定義できる - 弱依存境界まで行ける

最強条件: - 定理形で構造持続の数学に埋め込める

\section{運用上の原則}\label{ux904bux7528ux4e0aux306eux539fux5247}

原則1: 理論は薄く保つ。厚くするのは各ドメインの代理指標と実験である。

原則2: 100点の L は不要。80点でも今までない指標なら価値がある。ただし、80点の代理指標は「基準モデルより予測力がある」ことを示して初めて価値を持つ。「悪いことが増えると崩壊に近づく」という自明な予測だけでは不十分。

原則3: 最初から総合健全性を定義しない。維持条件は一つずつでよい。

原則4: 外れた代理指標は捨てる。理論を守るために代理指標を無理に残さない。

原則5: 理論の勝ち筋は、``すべてを最初から定理化すること''ではなく、``複数ドメインで同型の崩壊構造を示すこと''にある。

\section{LLM への適用例}\label{llm-ux3078ux306eux9069ux7528ux4f8b}

{\def\LTcaptype{table} % do not increment counter
\begin{longtable}[]{@{}
  >{\raggedright\arraybackslash}p{(\linewidth - 2\tabcolsep) * \real{0.5000}}
  >{\raggedright\arraybackslash}p{(\linewidth - 2\tabcolsep) * \real{0.5000}}@{}}
\toprule\noalign{}
\begin{minipage}[b]{\linewidth}\raggedright
項目
\end{minipage} & \begin{minipage}[b]{\linewidth}\raggedright
内容
\end{minipage} \\
\midrule\noalign{}
\endhead
\bottomrule\noalign{}
\endlastfoot
ドメイン & LLM \\
対象構造 & 論理一貫性を保った推論構造 \\
維持条件 & 論理一貫性維持 \\
崩壊条件 & 一貫した推論経路が閾値以下になること \\
反証条件 & (理論) 制約蓄積と無関係に崩壊が生じる / (代理指標) 矛盾密度が Y と無相関 \\
代理指標探索予算 & 主要代理指標 5個以内。全て外れたら理論の LLM 適用を再検討 \\
基準モデル & 「矛盾の総数が閾値を超えたら崩壊」（線形カウントモデル） \\
追加予測 & 制約の質による効き方の差（矛盾検出 \textgreater{} 規則 \textgreater{} 事実）、加法的累積構造 \\
有効選択肢空間 & 矛盾のない推論経路の集合 \\
制約集合 & 構造的矛盾、参照元衝突、指示衝突、記憶上書き \\
L̂ 代理指標 & 未解決矛盾密度、無効化推論分岐数 \\
μ & 推論余力、検索余力 \\
N\_eff\^{}(0) & 初期の独立な整合推論経路数 \\
ルート & 一般会話では Route C。明示制約タスクでは Route A または Route B \\
観測指標 & Y (論理一貫性維持率), z (矛盾密度), 検索成功率 \\
介入候補 & 矛盾解消（代謝）、時間的範囲づけ、記憶標識 \\
検証計画 & Stage B 実験（モデル横断の順序予測検証） \\
\end{longtable}
}

\section{本手順の目標}\label{ux672cux624bux9806ux306eux76eeux6a19}

本手順の目標は、各ドメインについて

維持条件 → 有効選択肢空間 → 制約削減 → L, μ, N\_eff\^{}(0) → 定理 or 境界 or 実証

へ写像する標準手順を与えることである。

これにより、分野ごとに別々の言葉で語られていた

\begin{itemize}
\tightlist
\item
  崩壊
\item
  回復力
\item
  劣化
\item
  代償不全
\item
  文脈腐敗
\item
  制度疲労
\item
  継続学習の失敗
\end{itemize}

を、同じ構造で比較・検証・介入できるようにする。

\section{結び}\label{ux7d50ux3073}

本補論が与えるのは、構造持続の最小形式を個別ドメインへ写像するための標準手順である。ここで重要なのは、すべてのドメインを最初から同じ強さで定理化することではない。むしろ、どこまでが等式として言え、どこからが境界になり、どこからが実証的主張になるかを明示したうえで、反証条件と基準モデル比較を伴う形で適用可能性を評価することにある。

したがって、本手順は完成した一般理論ではなく、複数ドメインで同型の崩壊構造を比較可能にするための作業仮説の枠組みである。その役割は、個別の実験や介入設計を散発的に並べるのではなく、共通の骨格のもとで積み上げられる形へ整えることにある。
\end{document}
